\section{Суффиксный массив и массив LCP}

На предыдущих лекциях мы изучали структуры, ориентированные на поиск фиксированных шаблонов в потоковом тексте. Теперь посмотрим на противоположную ситуацию: пусть у нас есть одна большая строка $s$ длины $n$, по которой мы хотим быстро отвечать на разнообразные запросы о её подстроках (например, сравнивать произвольные подстроки за $O(1)$, искать наибольшую общую подстроку или считать число различных подстрок). Основным инструментом для этого служит \textbf{суффиксный массив} в паре с \textbf{массивом $\lcp$}.

\subsection{Определения и базовые свойства}

Условимся обозначать через $s^i = s[i \dots n-1]$ суффикс строки $s$, начинающийся в позиции $i$ ($0 \le i < n$). Заметим, что все $n$ суффиксов строки $s$ имеют попарно различные длины (от $1$ до $n$), поэтому среди них не может быть двух одинаковых строк. Значит, все суффиксы можно однозначно упорядочить по возрастанию в стандартном лексикографическом порядке.

\begin{definition}
    \textbf{Суффиксным массивом} строки $s$ длины $n$ называется перестановка $(p_0, p_1, \dots, p_{n-1})$ чисел $\{0, 1, \dots, n-1\}$, задающая лексикографический порядок суффиксов строки $s$:
    \[
        s^{p_0} < s^{p_1} < \dots < s^{p_{n-1}}.
    \]
\end{definition}

Сам по себе суффиксный массив хранит лишь порядок суффиксов, но не говорит, насколько сильно похожи соседние суффиксы в этом порядке. Эту информацию несёт второй массив.

\begin{definition}
    Для произвольных строк $\alpha$ и $\beta$ обозначим через $\LCP(\alpha, \beta)$ (\textit{longest common prefix}) длину их наибольшего общего префикса. \textbf{Массивом $\lcp$} строки $s$ называется массив $(\lcp_0, \lcp_1, \dots, \lcp_{n-2})$ длины $n-1$, в котором $i$-й элемент равен длине наибольшего общего префикса соседних в суффиксном массиве суффиксов:
    \[
        \lcp_i = \LCP(s^{p_i}, s^{p_{i+1}}), \qquad 0 \le i \le n-2.
    \]
\end{definition}

\begin{example}
    Выпишем суффиксный массив и массив $\lcp$ для строки $s = \mathtt{ababba}$ длины $n = 6$:
    \begin{center}
    \begin{tabular}{@{}cccl@{}}
        \toprule
        Индекс $i$ & $p_i$ & $\lcp_i$ & Суффикс $s^{p_i}$ \\
        \midrule
        $0$ & $5$ & $1$ & $\mathtt{a}$ \\
        $1$ & $0$ & $2$ & $\mathtt{ababba}$ \\
        $2$ & $2$ & $0$ & $\mathtt{abba}$ \\
        $3$ & $4$ & $2$ & $\mathtt{ba}$ \\
        $4$ & $1$ & $1$ & $\mathtt{babba}$ \\
        $5$ & $3$ & --- & $\mathtt{bba}$ \\
        \bottomrule
    \end{tabular}
    \end{center}
\end{example}

Казалось бы, массив $\lcp$ хранит общие префиксы только для $(n-1)$ пар \emph{соседних} суффиксов, тогда как всего пар суффиксов $\binom{n}{2}$. Однако, благодаря тому что суффиксы отсортированы лексикографически, через массив $\lcp$ сразу выражается наибольший общий префикс вообще любой пары суффиксов.

\begin{lemma}[О наибольшем общем префиксе на отрезке суффиксного массива]\label{lem:lcp_min}
    Для любых двух индексов $0 \le l < r < n$ длина наибольшего общего префикса суффиксов $s^{p_l}$ и $s^{p_r}$ равна минимуму из значений $\lcp$ на полуинтервале $[l, r)$:
    \[
        \LCP(s^{p_l}, s^{p_r}) = \min_{l \le k < r} \lcp_k.
    \]
\end{lemma}

\begin{proof}
    Обозначим $m = \LCP(s^{p_l}, s^{p_r})$. Докажем равенство в две стороны.

    \textbf{1. Оценка $\min_{l \le k < r} \lcp_k \ge m$.} По определению, суффиксы $s^{p_l}$ и $s^{p_r}$ начинаются с одного и того же префикса $\alpha$ длины $m$. Так как в суффиксном массиве строки упорядочены лексикографически:
    \[
        s^{p_l} < s^{p_{l+1}} < \dots < s^{p_r},
    \]
    а крайние строки $s^{p_l}$ и $s^{p_r}$ начинаются с префикса $\alpha$, то и все промежуточные строки $s^{p_k}$ ($l \le k \le r$) обязаны начинаться с того же префикса $\alpha$ длины $m$. Следовательно, для любого $k \in [l, r-1]$ соседние суффиксы $s^{p_k}$ и $s^{p_{k+1}}$ совпадают хотя бы в первых $m$ символах, откуда $\lcp_k \ge m$.

    \textbf{2. Достижимость равенства.} По определению наибольшего общего префикса, в позиции $m$ суффиксы $s^{p_l}$ и $s^{p_r}$ различаются (либо суффикс $s^{p_l}$ закончился и имеет длину ровно $m$). Так как $s^{p_l} < s^{p_r}$, символ $s^{p_l}[m]$ строго меньше символа $s^{p_r}[m]$. Но при движении по индексам $k$ от $l$ к $r$ символ $s^{p_k}[m]$ в позиции $m$ не может оставаться неизменным: если в начале он равен $s^{p_l}[m]$, а в конце равен $s^{p_r}[m] > s^{p_l}[m]$, то найдётся хотя бы один шаг $k \in [l, r-1]$, на котором $s^{p_k}[m] < s^{p_{k+1}}[m]$. Для этого индекса $k$ суффиксы $s^{p_k}$ и $s^{p_{k+1}}$ совпадают в первых $m$ символах и различаются в $(m+1)$-м, то есть $\lcp_k = m$.
\end{proof}

\subsubsection{Сравнение произвольных подстрок за \texorpdfstring{$O(1)$}{O(1)}}

Лемма~\ref{lem:lcp_min} даёт удобный практический инструмент. Пусть мы построили суффиксный массив $p$, обратную к нему перестановку $\pos$ (где $\pos[i]$ "--- позиция суффикса $s^i$ в суффиксном массиве, то есть $\pos[p_k] = k$) и массив $\lcp$. Построим поверх массива $\lcp$ структуру данных \textbf{Sparse Table} (таблицу разреженных отрезков) за $O(n \log n)$ времени и памяти, которая позволяет находить минимум на любом отрезке массива за $O(1)$.

Тогда мы умеем за $O(1)$ находить $L = \LCP(s^i, s^j)$ для любых двух суффиксов $s^i$ и $s^j$: достаточно взять $\pos[i]$ и $\pos[j]$ (пусть для определённости $\pos[i] < \pos[j]$) и запросить минимум в Sparse Table на полуинтервале $[\pos[i], \pos[j])$.

Как с помощью этого сравнить две произвольные подстроки $A = s[l_1 \dots r_1]$ и $B = s[l_2 \dots r_2]$ за $O(1)$?
\begin{enumerate}
    \item Найдём $L = \LCP(s^{l_1}, s^{l_2})$ за $O(1)$.
    \item Если $L < \min(|A|, |B|)$, то первое различие между суффиксами $s^{l_1}$ и $s^{l_2}$ происходит строго внутри обеих подстрок $A$ и $B$. В этом случае достаточно сравнить символы $s[l_1 + L]$ и $s[l_2 + L]$ (или просто сравнить положения суффиксов $\pos[l_1]$ и $\pos[l_2]$).
    \item Если же $L \ge \min(|A|, |B|)$, то одна из подстрок является префиксом другой, и лексикографически меньше та подстрока, у которой меньше длина (а при $|A| = |B|$ подстроки совпадают).
\end{enumerate}
В отличие от полиномиальных хешей с бинарным поиском, которые сравнивают подстроки за $O(\log n)$ и имеют ненулевую вероятность коллизии, связка суффиксного массива и $\lcp$ отвечает на запрос сравнения за $O(1)$ и полностью детерминированно.

\subsection{Построение суффиксного массива за \texorpdfstring{$O(n \log n)$}{O(n log n)}}

Научимся строить суффиксный массив за время $O(n \log n)$. Для удобства сведём задачу сортировки суффиксов к задаче сортировки \textbf{циклических сдвигов} одинаковой длины.

Допишем в конец строки $s$ специальный символ-страж $\mathtt{\#}$ (например, символ с кодом $0$), который строго меньше любого символа алфавита $\Sigma$ и нигде в $s$ не встречается. Длина строки $s\mathtt{\#}$ станет равна $n + 1$ (далее для краткости будем обозначать новую длину снова через $n$).

\begin{proposition}
    Лексикографический порядок циклических сдвигов строки $s\mathtt{\#}$ совпадает с лексикографическим порядком её суффиксов.
\end{proposition}

\begin{proof}
    Рассмотрим два различных циклических сдвига строки $s\mathtt{\#}$, начинающихся в позициях $i \ne j$. Будем сравнивать их посимвольно слева направо. Поскольку символ $\mathtt{\#}$ входит в строку $s\mathtt{\#}$ ровно один раз, в двух различных сдвигах он стоит на разных расстояниях от начала. В тот момент, когда в одном из сдвигов встретится символ $\mathtt{\#}$, в другом сдвиге на этой же позиции будет стоять обычный символ из $\Sigma$, который строго больше $\mathtt{\#}$. Значит, первое несовпадение между циклическими сдвигами гарантированно произойдёт не позже, чем на позиции символа $\mathtt{\#}$. Но до символа $\mathtt{\#}$ циклический сдвиг в точности совпадает с соответствующим суффиксом исходной строки $s$.
\end{proof}

После того как мы отсортируем все циклические сдвиги строки $s\mathtt{\#}$, самый первый элемент в отсортированном порядке будет начинаться с самого символа $\mathtt{\#}$ (нулевой суффикс); отбросив его, мы получим искомый суффиксный массив строки $s$.

\subsubsection{Итеративное удвоение и цифровая сортировка}

Будем итеративно сортировать циклические подстроки длины $2^k$ при $k = 0$, $1$, $\dots$, $\lceil \log_2 n \rceil$. Под циклической подстрокой длины $\ell$ с началом в позиции $i$ понимается строка
\[
    s[i] \, s[(i+1) \bmod n] \dots s[(i + \ell - 1) \bmod n].
\]
На $k$-й итерации мы хотим поддерживать два массива длины $n$:
\begin{itemize}
    \item перестановку $(p_0, p_1, \dots, p_{n-1})$, упорядочивающую циклические подстроки длины $2^k$ по неубыванию;
    \item массив \textbf{классов эквивалентности} $(c_0, c_1, \dots, c_{n-1})$, где $c_i \in \{0, 1, \dots, n-1\}$ "--- номер класса циклической подстроки длины $2^k$, начинающейся в позиции $i$. При этом равные подстроки получают одинаковый номер класса, а лексикографически меньшая подстрока получает строго меньший номер класса.
\end{itemize}

\textbf{База ($k = 0$, длина $2^0 = 1$).} Циклические подстроки длины $1$ "--- это отдельные символы $s[i]$. Мы можем отсортировать их обычной сортировкой подсчётом (Counting Sort) за время $O(n + |\Sigma|)$, после чего пройтись по отсортированному массиву $p$ и выставить классы эквивалентности $c_{p_i}$ (увеличивая номер класса на $1$ всякий раз, когда $s[p_i] \ne s[p_{i-1}]$).

\textbf{Переход ($k \to k + 1$, от длины $2^k$ к длине $2^{k+1}$).} Пусть мы уже умеем сравнивать циклические подстроки длины $2^k$ за $O(1)$ с помощью массива классов $c$. Любая циклическая подстрока длины $2^{k+1}$ с началом в позиции $i$ распадается на две половинки длины $2^k$: левая половинка начинается в позиции $i$, а правая "--- в позиции $(i + 2^k) \bmod n$.

\begin{center}
\begin{tikzpicture}[scale=1.0]
    \draw[thick, fill=NavyBlue!10] (0, 0) rectangle (5.4, 0.75) node[pos=0.5, font=\footnotesize] {левая половина: класс $c_i$};
    \draw[thick, fill=OliveGreen!15] (5.4, 0) rectangle (12.0, 0.75) node[pos=0.5, font=\footnotesize] {правая половина: класс $c_{(i + 2^k) \bmod n}$};
    \draw[decorate, decoration={brace, amplitude=6pt}, thick] (0, 0.9) -- (12.0, 0.9) node[midway, above=6pt, font=\small] {циклическая подстрока длины $2^{k+1}$};
    \node[below, font=\footnotesize] at (0, 0) {$i$};
    \node[below, font=\footnotesize] at (5.4, 0) {$(i + 2^k) \bmod n$};
\end{tikzpicture}
\end{center}

Следовательно, сравнение двух подстрок длины $2^{k+1}$ с началами в $i$ и $j$ равносильно лексикографическому сравнению двух пар целых чисел:
\[
    \bigl(c_i,\; c_{(i + 2^k) \bmod n}\bigr) \quad \text{и} \quad \bigl(c_j,\; c_{(j + 2^k) \bmod n}\bigr).
\]
Если мы отсортируем эти $n$ пар обычным \texttt{std::sort}, то потратим $O(n \log n)$ на одну итерацию и $O(n \log^2 n)$ суммарно. Но обе компоненты каждой пары "--- целые числа из диапазона от $0$ до $n-1$. Поэтому $n$ пар можно отсортировать за линейное время $O(n)$ с помощью \textbf{цифровой сортировки (Radix Sort)}: сначала отсортировать по младшему разряду (второму элементу пары), а затем устойчиво отсортировать по старшему разряду (первому элементу пары).

Ключевое наблюдение состоит в том, что \textbf{первая сортировка (по второму элементу пары) вообще не требует работы}. В самом деле, второй элемент пары для позиции $i$ "--- это класс подстроки длины $2^k$, начинающейся в $(i + 2^k) \bmod n$. Но в массиве $p$ с предыдущего шага все подстроки длины $2^k$ \emph{уже} отсортированы. Значит, если мы просто сдвинем все индексы из массива $p$ на $2^k$ позиций влево по модулю $n$:
\[
    p'_i = (p_i - 2^k + n) \bmod n,
\]
то в массиве $p'$ индексы будут идти в точности в порядке неубывания их вторых половинок $c_{(p'_i + 2^k) \bmod n} = c_{p_i}$.

Остаётся выполнить всего одну устойчивую сортировку подсчётом массива $p'$ по значениям первых элементов пар $c_{p'_i}$, получая новый массив $p$, и за один линейный проход пересчитать новый массив классов эквивалентности $c'$. Так как на каждой из $\lceil \log_2 n \rceil$ итераций мы тратим $O(n)$ времени, суммарная асимптотика алгоритма составляет $O(n \log n)$.

\begin{lstlisting}[caption={Построение суффиксного массива за $O(n \log n)$}]
vector<int> build_suffix_array(string s) {
    s += char(0); // добавляем символ-страж '\#' с кодом 0
    int n = s.size();
    const int ALPHA = 256;

    vector<int> p(n), c(n), cnt(max(ALPHA, n), 0);

    // База $k = 0$: сортировка подсчётом одиночных символов
    for (int i = 0; i < n; ++i) {
        ++cnt[(unsigned char)s[i]];
    }
    for (int i = 1; i < ALPHA; ++i) {
        cnt[i] += cnt[i - 1];
    }
    for (int i = n - 1; i >= 0; --i) {
        p[--cnt[(unsigned char)s[i]]] = i;
    }

    c[p[0]] = 0;
    int classes = 1;
    for (int i = 1; i < n; ++i) {
        if (s[p[i]] != s[p[i - 1]]) {
            ++classes;
        }
        c[p[i]] = classes - 1;
    }

    // Шаги удвоения длины: от $2^k$ к $2^{k+1}$
    vector<int> pn(n), cn(n);
    for (int len = 1; len < n; len <<= 1) {
        // 1. Сортировка по второй половине за счёт циклического сдвига
        for (int i = 0; i < n; ++i) {
            pn[i] = p[i] - len;
            if (pn[i] < 0) pn[i] += n;
        }

        // 2. Устойчивая сортировка подсчётом по первой половине
        fill(cnt.begin(), cnt.begin() + classes, 0);
        for (int i = 0; i < n; ++i) {
            ++cnt[c[pn[i]]];
        }
        for (int i = 1; i < classes; ++i) {
            cnt[i] += cnt[i - 1];
        }
        for (int i = n - 1; i >= 0; --i) {
            p[--cnt[c[pn[i]]]] = pn[i];
        }

        // 3. Пересчёт классов эквивалентности для длины 2 * len
        cn[p[0]] = 0;
        classes = 1;
        for (int i = 1; i < n; ++i) {
            pair<int, int> cur = {c[p[i]], c[(p[i] + len) % n]};
            pair<int, int> prev = {c[p[i - 1]], c[(p[i - 1] + len) % n]};
            if (cur != prev) {
                ++classes;
            }
            cn[p[i]] = classes - 1;
        }
        c.swap(cn);
        if (classes == n) break; // все суффиксы уже попарно различны
    }

    // Удаляем нулевой циклический сдвиг, начинающийся с символа-стража
    p.erase(p.begin());
    return p;
}
\end{lstlisting}

\begin{note}[Применение в обучении больших языковых моделей]
    Итеративный алгоритм построения суффиксного массива неожиданно пригодился в современном машинном обучении. В 2022 году исследователи Google опубликовали статью \textit{<<Deduplicating Training Data Makes Language Models Better>>} (Lee et al., ACL 2022). При обучении LLM на гигантских корпусах из интернета (сотни гигабайт и терабайты текста) в данных неизбежно встречаются тысячи одинаковых фрагментов "--- шаблонные куски сайтов, перепечатки новостей, скопированные абзацы. Из-за них модель переобучается и начинает дословно <<зубрить>> дубликаты.

    Как найти все совпадающие подстроки длины хотя бы $64$ токена в корпусе из ста миллиардов токенов? Строить полный суффиксный массив до конца слишком долго, но этого и не нужно: достаточно запустить описанный выше алгоритм удвоения ровно на $6$ итераций ($2^6 = 64$). После шестой итерации все подстроки длины $64$, попавшие в один и тот же класс эквивалентности $c_i$, гарантированно совпадают, и их можно отфильтровать за линейный проход.
\end{note}

\subsection{Алгоритм Касаи построения массива LCP за \texorpdfstring{$O(n)$}{O(n)}}

Пусть суффиксный массив $p$ уже построен. Как теперь найти массив $\lcp$? Если сравнивать соседние суффиксы $s^{p_i}$ и $s^{p_{i+1}}$ наивно посимвольно, на строке из одних букв $\mathtt{a}$ мы получим квадратичное время $\Theta(n^2)$.

В 2001 году Тору Касаи с соавторами заметил, что массив $\lcp$ можно вычислить за линейное время $O(n)$, если вычислять его элементы не в порядке индексов суффиксного массива $0, 1, \dots, n-2$, а \textbf{в порядке убывания длины суффиксов} "--- от самого длинного суффикса $s^0$ к самому короткому суффиксу $s^{n-1}$.

Обозначим через $\pos[i]$ позицию суффикса $s^i$ в суффиксном массиве (так что $p_{\pos[i]} = i$).

\begin{lemma}[Ключевая лемма алгоритма Касаи]\label{lem:kasai}
    Пусть для некоторого $0 \le i < n - 1$ суффикс $s^i$ не является последним в суффиксном массиве ($\pos[i] < n - 1$), и пусть $\lcp_{\pos[i]} = k > 0$. Тогда для следующего суффикса $s^{i+1}$ (если $\pos[i+1] < n - 1$) справедливо неравенство
    \[
        \lcp_{\pos[i+1]} \ge k - 1.
    \]
\end{lemma}

\begin{proof}
    Пусть $r = \pos[i]$ и пусть $j = p_{r+1}$ "--- индекс суффикса, который стоит в суффиксном массиве непосредственно следом за $s^i$. По условию, $\LCP(s^i, s^j) = \lcp_r = k \ge 1$, и при этом $s^i < s^j$.

    Так как $k \ge 1$, первые $k$ символов у суффиксов $s^i$ и $s^j$ совпадают:
    \[
        s[i \dots i + k - 1] = s[j \dots j + k - 1].
    \]
    Отбросим первый символ у обоих суффиксов $s^i$ и $s^j$ "--- мы получим суффиксы $s^{i+1}$ и $s^{j+1}$. Что можно сказать про эту пару суффиксов?
    \begin{enumerate}
        \item Поскольку у $s^i$ и $s^j$ совпадали первые $k$ символов, после удаления первого символа у $s^{i+1}$ и $s^{j+1}$ совпадают первые $k - 1$ символов, то есть $\LCP(s^{i+1}, s^{j+1}) = k - 1$.
        \item Поскольку $s^i < s^j$ и их первые символы были одинаковы ($s[i] = s[j]$), лексикографический порядок при удалении общей первой буквы сохраняется: $s^{i+1} < s^{j+1}$.
    \end{enumerate}

    \begin{center}
    \begin{tikzpicture}[scale=1.0]
        % Suffix array column representation
        \node[font=\small\bfseries] at (-1.6, 2.2) {Позиция в SA};
        \node[font=\small\bfseries] at (2.8, 2.2) {Суффикс};

        % s^i and s^j
        \node[font=\small] at (-1.6, 1.4) {$\pos[i]$};
        \draw[thick, fill=NavyBlue!15] (0, 1.15) rectangle (0.7, 1.65) node[pos=0.5, font=\footnotesize] {$c$};
        \draw[thick, fill=NavyBlue!15] (0.7, 1.15) rectangle (3.9, 1.65) node[pos=0.5, font=\footnotesize] {$\alpha$ (длина $k-1$)};
        \draw[thick] (3.9, 1.15) rectangle (5.7, 1.65) node[pos=0.5, font=\footnotesize] {$\dots$};
        \node[right, font=\small] at (5.9, 1.4) {$= s^i$};

        \node[font=\small] at (-1.6, 0.7) {$\pos[i]+1$};
        \draw[thick, fill=NavyBlue!15] (0, 0.45) rectangle (0.7, 0.95) node[pos=0.5, font=\footnotesize] {$c$};
        \draw[thick, fill=NavyBlue!15] (0.7, 0.45) rectangle (3.9, 0.95) node[pos=0.5, font=\footnotesize] {$\alpha$ (длина $k-1$)};
        \draw[thick] (3.9, 0.45) rectangle (6.0, 0.95) node[pos=0.5, font=\footnotesize] {$\dots$};
        \node[right, font=\small] at (6.2, 0.7) {$= s^j$};

        % s^{i+1} and s^{j+1}
        \node[font=\small] at (-1.6, -0.5) {$\pos[i+1]$};
        \draw[thick, fill=OliveGreen!15] (0, -0.75) rectangle (3.2, -0.25) node[pos=0.5, font=\footnotesize] {$\alpha$ (длина $k-1$)};
        \draw[thick] (3.2, -0.75) rectangle (5.0, -0.25) node[pos=0.5, font=\footnotesize] {$\dots$};
        \node[right, font=\small] at (5.2, -0.5) {$= s^{i+1}$};

        \node[font=\small] at (-1.6, -1.2) {$\pos[i+1]+1$};
        \draw[thick, fill=OliveGreen!15] (0, -1.45) rectangle (3.2, -0.95) node[pos=0.5, font=\footnotesize] {$\alpha$ (длина $k-1$)};
        \draw[thick] (3.2, -1.45) rectangle (5.2, -0.95) node[pos=0.5, font=\footnotesize] {$\dots$};
        \node[right, font=\small] at (5.4, -1.2) {$= s^{p_{\pos[i+1]+1}}$};

        \node[font=\small] at (-1.6, -1.7) {$\vdots$};
        \node[font=\small] at (2.4, -1.7) {$\vdots$};

        \node[font=\small] at (-1.6, -2.3) {$\pos[j+1]$};
        \draw[thick, fill=OliveGreen!15] (0, -2.55) rectangle (3.2, -2.05) node[pos=0.5, font=\footnotesize] {$\alpha$ (длина $k-1$)};
        \draw[thick] (3.2, -2.55) rectangle (5.3, -2.05) node[pos=0.5, font=\footnotesize] {$\dots$};
        \node[right, font=\small] at (5.5, -2.3) {$= s^{j+1}$};
    \end{tikzpicture}
    \end{center}

    Из неравенства $s^{i+1} < s^{j+1}$ следует, что в суффиксном массиве строка $s^{j+1}$ располагается строго ниже (правее) строки $s^{i+1}$, то есть $\pos[i+1] < \pos[j+1]$. Применим теперь лемму~\ref{lem:lcp_min} о минимуме на отрезке к паре индексов $l = \pos[i+1]$ и $r = \pos[j+1]$:
    \[
        k - 1 = \LCP(s^{i+1}, s^{j+1}) = \min_{\pos[i+1] \le t < \pos[j+1]} \lcp_t \le \lcp_{\pos[i+1]},
    \]
    так как $\lcp_{\pos[i+1]}$ входит в множество чисел, по которым берётся минимум.
\end{proof}

Лемма~\ref{lem:kasai} сразу подсказывает алгоритм: будем перебирать $i$ от $0$ до $n-1$, поддерживая текущее значение $k = \lcp_{\pos[i]}$. При переходе от $i$ к $i + 1$ мы знаем, что первые $\max(0, k - 1)$ символов у суффикса $s^{i+1}$ и его соседа справа в суффиксном массиве заведомо совпадают. Поэтому их не нужно проверять заново "--- мы просто уменьшаем $k$ на $1$ (если $k > 0$) и дошагиваем посимвольным сравнением, начиная сразу с позиции $k$.

\begin{lstlisting}[caption={Алгоритм Касаи построения массива $\lcp$ за $O(n)$}]
vector<int> build_lcp(const string& s, const vector<int>& p) {
    int n = s.size();
    vector<int> pos(n);
    for (int i = 0; i < n; ++i) {
        pos[p[i]] = i;
    }

    vector<int> lcp(n - 1, 0);
    int k = 0;
    for (int i = 0; i < n; ++i) {
        if (pos[i] == n - 1) {
            k = 0;
            continue;
        }
        int j = p[pos[i] + 1]; // сосед суффикса $s^i$ справа в суффиксном массиве
        while (i + k < n && j + k < n && s[i + k] == s[j + k]) {
            ++k;
        }
        lcp[pos[i]] = k;
        if (k > 0) {
            --k;
        }
    }
    return lcp;
}
\end{lstlisting}

\begin{theorem}
    Алгоритм Касаи вычисляет массив $\lcp$ за время $O(n)$.
\end{theorem}

\begin{proof}
    Оценим количество итераций внутреннего цикла \texttt{while} за всё время работы алгоритма. На каждой успешной итерации цикла \texttt{while} переменная $k$ увеличивается ровно на $1$. В то же время значение $k$ никогда не превосходит $n$, в начале работы равно $0$, а уменьшаться может только в конце шага внешнего цикла \texttt{for} в строчке \texttt{if (k > 0) -\{\}-k;} "--- то есть не более чем на $1$ за каждую из $n$ итераций внешнего цикла.

    Раз переменная $k \le n$ за всё время работы уменьшается суммарно не более чем на $n$, она не может и увеличиться более чем $2n$ раз! Значит, внутренний цикл \texttt{while} выполнит суммарно не более $2n$ проверок, и общее время работы алгоритма равно $O(n)$.
\end{proof}

\begin{problem}[Количество различных подстрок]
    Дана строка $s$ длины $n$. Найдите количество её попарно различных непустых подстрок за время $O(n \log n)$.
\end{problem}

\begin{solution}
    Любая подстрока строки $s$ "--- это префикс некоторого её суффикса. Будем перебирать суффиксы в порядке суффиксного массива: $s^{p_0}, s^{p_1}, \dots, s^{p_{n-1}}$. Первый суффикс $s^{p_0}$ имеет длину $n - p_0$ и даёт нам $n - p_0$ различных префиксов. Каждый следующий суффикс $s^{p_i}$ ($i \ge 1$) имеет длину $n - p_i$, но ровно $\lcp_{i-1}$ его кратчайших префиксов уже встречались как префиксы предыдущего суффикса $s^{p_{i-1}}$ (а по лемме~\ref{lem:lcp_min} с более ранними суффиксами общий префикс не может быть длиннее $\lcp_{i-1}$). Следовательно, суффикс $s^{p_i}$ добавляет ровно $(n - p_i) - \lcp_{i-1}$ новых подстрок, а общее число различных подстрок равно
    \[
        \sum_{i=0}^{n-1} (n - p_i) - \sum_{i=0}^{n-2} \lcp_i = \frac{n(n+1)}{2} - \sum_{i=0}^{n-2} \lcp_i.
    \]
\end{solution}
